\setchapterpreamble[u]{\margintoc}
\chapter{Expectation, Variance, and Covariance}
\labch{expectation}

\sources{Bertsekas and Tsitsiklis, \emph{Introduction to Probability}, 2nd ed.\
\cite{bertsekas2008probability}, Chapters~2--4; MIT 6.041 \cite{ocw6041},
Lectures~5--12.}

Distributions are unwieldy objects. Expectation reduces one to a single number,
variance to a second, and covariance extends the idea to pairs --- producing, in
the end, the covariance matrix already met in \refch{learning-from-data}.

\section{Expectation}
\labsec{expectation-def}

\begin{definition}
The \emph{expectation} of \(X\) is
\[
	\E[X]=\sum_x x\,p_X(x)
	\qquad\text{or}\qquad
	\E[X]=\int x\,f_X(x)\,dx,
\]
according as \(X\) is discrete or continuous, when the sum or integral converges
absolutely.
\end{definition}

\begin{theorem}[Linearity]
For random variables \(X,Y\) on the same space and constants \(a,b\),
\[
	\E[aX+bY]=a\E[X]+b\E[Y].
\]
No independence is required.
\end{theorem}

\marginnote{Linearity without independence is the workhorse of the subject. It is
what makes the mini-batch gradient of \refsec{sgd} unbiased even though the
examples in a batch share a model and a step.}

\section{Expectations of functions}
\labsec{lotus}

The expectation of \(g(X)\) needs no new distribution:
\[
	\E[g(X)]=\sum_x g(x)p_X(x)
	\qquad\text{or}\qquad
	\E[g(X)]=\int g(x)f_X(x)\,dx .
\]
Contrast \refsec{transformations}: obtaining the \emph{density} of \(g(X)\)
requires a Jacobian factor, obtaining its \emph{expectation} does not.

\begin{remark}
Except for affine \(g\), \(\E[g(X)]\neq g(\E[X])\). For convex \(g\) the
inequality has a direction --- Jensen's inequality gives
\(\E[g(X)]\ge g(\E[X])\) --- which is why an average of losses is not the loss of
an average.
\end{remark}

\section{Variance}
\labsec{variance}

\begin{definition}
\(\Var(X)=\E\bigl[(X-\E[X])^{2}\bigr]\), and the standard deviation is its
square root.
\end{definition}

Expanding the square gives the computational form
\(\Var(X)=\E[X^{2}]-\E[X]^{2}\), and scaling behaves quadratically:
\(\Var(aX+b)=a^{2}\Var(X)\). For \emph{independent} \(X\) and \(Y\),
\(\Var(X+Y)=\Var(X)+\Var(Y)\).

\begin{example}
For \(X\sim\)~Bernoulli\((p)\), \(\E[X]=p\) and \(\E[X^{2}]=p\), so
\(\Var(X)=p-p^{2}=p(1-p)\), largest at \(p=\tfrac12\) and zero at the extremes.
\end{example}

\begin{proposition}
If \(X_1,\ldots,X_n\) are i.i.d.\ with mean \(\mu\) and variance
\(\sigma^{2}\), the sample mean \(\bar X_n=\frac1n\sum_iX_i\) satisfies
\[
	\E[\bar X_n]=\mu,
	\qquad
	\Var(\bar X_n)=\frac{\sigma^{2}}{n}.
\]
\end{proposition}

\marginnote{The \(1/n\) --- and therefore the \(1/\sqrt n\) in the standard
deviation --- is the reason quadrupling a batch halves the gradient noise, and
the reason the returns diminish so quickly.}

\section{Covariance and correlation}
\labsec{covariance-def}

\begin{definition}
\(\Cov(X,Y)=\E\bigl[(X-\E[X])(Y-\E[Y])\bigr]\), and the correlation is
\(\rho=\Cov(X,Y)/\sqrt{\Var(X)\Var(Y)}\in[-1,1]\).
\end{definition}

Independence implies zero covariance; the converse is false. A standard
counterexample is \(X\) symmetric about zero with \(Y=X^{2}\): the two are
completely dependent, yet \(\Cov(X,Y)=\E[X^{3}]=0\).

\begin{remark}
Covariance detects linear association only. Since it is exactly what principal
component analysis in \refsec{pca} consumes, that method sees linear structure
and nothing else --- a limitation of the tool, not of the data.
\end{remark}

\section{The covariance matrix}
\labsec{covariance-matrix}

For a random vector \(\vect{X}\in\R^{d}\) with mean \(\vgreek{\mu}\),
\[
	\Sigma=\E\bigl[(\vect{X}-\vgreek{\mu})(\vect{X}-\vgreek{\mu})\tp\bigr]\in\R^{d\times d}.
\]

\begin{proposition}
\(\Sigma\) is symmetric positive semidefinite, and for any fixed
\(\vect{a}\in\R^{d}\),
\[
	\Var(\vect{a}\tp\vect{X})=\vect{a}\tp\Sigma\vect{a}.
\]
\end{proposition}

\begin{proof}
Symmetry is immediate from the definition. For the second claim, expand
\(\Var(\vect{a}\tp\vect{X})
=\E\bigl[(\vect{a}\tp(\vect{X}-\vgreek{\mu}))^{2}\bigr]
=\vect{a}\tp\E[(\vect{X}-\vgreek{\mu})(\vect{X}-\vgreek{\mu})\tp]\vect{a}\).
Since a variance is non-negative, the quadratic form is non-negative for every
\(\vect{a}\), which is positive semidefiniteness.
\end{proof}

\marginnote{So the population covariance matrix and the sample covariance matrix
\(S\) of \refsec{covariance} are the same object, one exact and one estimated.
The spectral theorem applies to both, which is what principal components rely
on.}

\section{Conditional expectation}
\labsec{conditional-expectation}

\(\E[Y\mid X=x]\) is the mean of the conditional distribution, and
\(\E[Y\mid X]\) is the random variable taking that value when \(X=x\). The tower
property states
\[
	\E\bigl[\E[Y\mid X]\bigr]=\E[Y].
\]

\begin{proposition}
Among all functions \(g\), the one minimising \(\E\bigl[(Y-g(X))^{2}\bigr]\) is
\(g(X)=\E[Y\mid X]\).
\end{proposition}

\begin{remark}
This identifies what a regression model trained on squared error is trying to
be: the conditional mean of the target given the input. Everything the model
gets wrong is either approximation error --- the family cannot express that
conditional mean --- or estimation error from finite data.
\end{remark}
